\chapter{Signature Games, Theorem Chain, and Paper Proof}
\label{ch:games}

\section{Generic ROM Signature Interface}

The file \code{Sig\_ROM.ec} defines the generic games.  Important declarations
include \code{Oracle}, \code{POracle}, \code{Scheme}, \code{SignOracle},
\code{Adversary}, \code{NMA\_Adversary}, \code{EUF\_CMA}, and \code{UF\_NMA}.
These definitions give the proof a stable target: \haetae{} security means
security in this game interface.

The distinction between EUF-CMA and UF-NMA is proof-theoretic.  The chosen
message game gives the adversary an adaptive signing oracle.  The no-message
game removes that oracle and is therefore closer to the algebraic extraction
surface.  Much of the theorem chain is devoted to justifying that signing-oracle
access can be simulated and internalized without changing the adversary's view
except for explicitly bounded losses.  In paper notation this is often a single
sentence; in the checked development it becomes a sequence of modules and
relational lemmas.

\section{Game-Hopping Calculus}

Let \(G_0,\ldots,G_n\) be a sequence of games with the same Boolean success
event.  The standard telescoping inequality is
\[
  \left|\Pr[G_0\Rightarrow 1]-\Pr[G_n\Rightarrow 1]\right|
  \le
  \sum_{i=0}^{n-1}
  \left|\Pr[G_i\Rightarrow 1]-\Pr[G_{i+1}\Rightarrow 1]\right|.
\]
Each adjacent term must be justified by either perfect equivalence, a
fundamental-lemma bad-event bound, or a reduction to a hardness assumption.  The
HAETAE proof is long because the adjacent games are not cosmetic.  They expose
oracle logs, signing transcripts, sampler sites, rejection events, and exact
sampling distributions at different stages.

This calculus also explains why loss ledgers are delicate.  A loss should appear
exactly once: either as a local bad-event probability, as a statistical-distance
term, or as part of a reduction advantage.  Double-counting makes the theorem
weaker and conceptually inaccurate; omitting a term makes it false.

\section{Pen-and-Paper EUF-CMA Game}

Let \(\Pi=(\mathsf{KeyGen},\mathsf{Sign},\mathsf{Verify})\) be the modeled
\haetae{} signature scheme.  The \eufcma{} experiment is:

\begin{enumerate}
\item Sample \((pk,sk)\leftarrow \mathsf{KeyGen}()\).
\item Give \(pk\) and random-oracle access to an adversary \(A\).
\item Answer each signing query \((m,ctx)\) with
\(\sigma\leftarrow\mathsf{Sign}(sk,m,ctx)\), and record \((m,ctx)\).
\item Let \(A\) output \((m^*,ctx^*,\sigma^*)\).
\item The adversary wins if
\(\mathsf{Verify}(pk,m^*,ctx^*,\sigma^*)=1\) and \((m^*,ctx^*)\) was not a
previous signing query.
\end{enumerate}

The advantage is the probability of this event over key generation, oracle
randomness, signing randomness, and adversarial randomness.

The resource parameters \(q_H\) and \(q_S\) are part of the paper-style loss
ledger and of the budgeted wrapper games, not syntactic caps on the source
\code{Sig\_ROM.EUF\_CMA} game.  In a concrete proof, these parameters enter
collision bounds, freshness bounds, and wrapper lemmas.  A theorem checked only
for fixed wrapper constants \(q_H,q_S\) should not be cited as a symbolic-query
or arbitrary-query theorem without a separate parameterized development.

\paragraph{Schematic paper-style ledger.}
A fully parameterized paper theorem would have the form
\[
\Adv^{\eufcma}_{\haetae}(A)
\le
\sum_i \Adv_i(B_i) + \mathsf{Loss}_{\mathrm{ROM}}(q_H,q_S)
+ \mathsf{Loss}_{\mathrm{local}}(q_H,q_S),
\]
where \(\mathsf{Loss}_{\mathrm{local}}\) denotes intermediate rejection,
sampling, and lifting costs before they are composed into a final implication.
The dissertation does not claim that this symbolic-query statement has been
checked as written.  The checked counted-ROM implication fixes \(q_H=q_S=2\) and
uses the concrete bound \code{haetae\_euf\_counted\_rom\_bound}; in that final
bound, rejection and exact-sampling obligations are discharged inside the hop
premises rather than added again as separate final summands.

\begin{proof}[Pen-and-paper proof sketch]
The proof proceeds by a sequence of games.  First, replace real signing by a
signing simulator while preserving the adversary-visible distribution.  Second,
expose internal signing transcripts and erase components not visible to the
adversary.  Third, move from chosen-message access to an NMA-style game by
embedding signatures into simulated transcript generation.  Fourth, program the
random oracle at challenge sites and account for the event that the adversary
prequeried a programmed point.  Fifth, pass through the rejection-sampling and
paper-simulation sampler surface, recording the checked rejection and lifting
obligations.  Sixth, replace the paper/structural sampler view by the exact
hyperball sampler used in the formal model, then apply special soundness and the
bimodal-to-Module-SIS bridge to turn a forgery into a lattice-assumption
adversary.  Summing the statistical distances and bad-event probabilities gives
the schematic inequality above; the machine-checked counted theorem uses the
fixed-ledger composition and grouping described later in the dissertation.
\end{proof}

The proof sketch is intentionally schematic.  The checked proof strengthens it
by replacing each informal distributional claim with an EasyCrypt theorem.  For
example, ``preserving the adversary-visible distribution'' becomes a relational
equivalence between games; ``account for the event that the adversary prequeried
a programmed point'' becomes a logged-ROM bad-event probability bound; and
``summing the losses'' becomes arithmetic over named bound operations in
\code{HAETAE\_Reductions.ec}.

\section{Machine-Checked Theorem Chain}

\code{HAETAE\_Security.ec} contains the top-level composition chain.  Representative
lemmas include \code{correctness\_perfect}, \code{euf\_cma\_security},
\code{euf\_cma\_security\_with\_hop\_interfaces}, and later bridge lemmas that
connect ROM transcript games, counted bimodal reductions, paper-simulation NMA
games, rejection sampling, and exact hyperball sampling.

The proof follows the paper argument, but each paper step is made explicit:

\begin{longtable}{p{0.33\textwidth}p{0.56\textwidth}}
\toprule
\textbf{Paper step} & \textbf{EasyCrypt realization} \\
\midrule
Real EUF-CMA game & \code{Sig\_ROM.EUF\_CMA} and \code{HAETAE\_Security.euf\_cma\_security}. \\
Signing simulation & \code{EUF\_CMA\_SimulatedSign} and \code{RealSigningSimulator}. \\
Transcript exposure & \code{EUF\_CMA\_TranscriptSimulatedSign}, \code{EUF\_CMA\_InternalTranscriptSign}. \\
ROM-internal transcript game & \code{EUF\_CMA\_ROMInternalTranscriptSign}. \\
NMA reduction surface & \code{ROMInternalTranscriptAsNMA} and budgeted variants. \\
Paper sampling to exact sampling & \code{ROSigningAttemptPaperSimSampler}, \code{ROExactHyperballPaperSimSampler}. \\
Final concrete bound & \code{HAETAE\_Reductions.ec} grouped bound lemmas. \\
\bottomrule
\end{longtable}

\section{Detailed EasyCrypt Code Analysis}

The design of \code{Sig\_ROM.ec} is intentionally parametric.  The module type
\code{Scheme(O : POracle)} prevents the proof from assuming a concrete hash
implementation in the generic security game.  The adversary receives only the
interfaces exposed by \code{POracle} and \code{SignOracle}; therefore, any
state used internally by the simulator must be hidden unless explicitly exposed
by a later game.

The \haetae{}-specific theorem file does not redefine security.  Instead, it
instantiates and bounds the generic security game.  This is a crucial design
choice.  If \code{HAETAE\_Security.ec} had its own custom forgery game, the
proof would also need a separate theorem showing equivalence with the standard
signature notion.  By reusing \code{Sig\_ROM.ec}, the development makes the
security target auditable.

The long chain of lemmas in \code{HAETAE\_Security.ec} should be read as a
machine-checked proof outline.  Lemmas with names containing
\code{from\_rom\_transcript}, \code{from\_structural\_nma},
\code{from\_paper\_sim\_nma}, and \code{from\_exact\_hyperball} correspond to
successive paper-proof transformations.  This naming discipline is valuable:
it lets a reviewer locate the exact checked boundary for each cryptographic
claim.
